\begin{abstracts}

An always-on robot can learn the social rhythm of a room over several days and still behave as it did on the first morning. The always-on cognitive architecture developed for the iCub segments a continuous audiovisual stream into episodes and recognises, without supervision, when a shared space is lively and when it is calm, yet its behaviour never draws on that knowledge. Its authors name what is missing: an action-selection loop that turns recognition into action, and a value system that defines what matters to the robot.

This thesis builds the loop, an embodied behaviour architecture that replaces the original behaviour module, and installs one value on it. The architecture tracks people as individuals and chooses whom to attend to from salience and learned affinity, while the robot's internal state sets the priority of what it does and says, in person and over Telegram. The value is a synthetic orexigenic drive: an energy reserve that drains with time and with every social act, and that only the people nearby can refill. Energy maintenance, which robots usually leave to a docking station or an operator, becomes part of the social exchange. Each regulatory outcome also updates a per-person affinity, so the robot remembers who fed it and favours them in later approaches.

An eight-day office deployment with twelve participants tested the loop. Deficit detection and behavioural redirection matched the design. Deficit was associated with more feeding in both phases; meals grew with deficit, and affinity changed approach eligibility. Recovery depended on caregiver presence; remote requests were rarely answered. Mean hunger ratings were below the midpoint for naturalness and genuineness and above it for helping motivation, without significant midpoint differences after correction. Helping motivation exceeded naturalness in an exploratory paired comparison. Written accounts identified repetitive, poorly timed requests and unchanging non-verbal expression. Regulatory effectiveness and credibility diverged. Both require evaluation, including how wording, timing and non-verbal expression shape credibility.

\end{abstracts}
